\documentclass[11pt]{article}

% --- Core style mirrors teaching-placement lesson packets (article + simple packages)
\usepackage[utf8]{inputenc}
\usepackage[T1]{fontenc}
\usepackage[margin=0.6in]{geometry}
\usepackage{hyperref}
\usepackage{enumitem}

% --- Helpful formatting for a full multi-session lesson
\usepackage{booktabs}
\usepackage{tabularx}
\usepackage{longtable}
\usepackage{xcolor}
\usepackage{array}

\title{MagicSchool Choose-Your-Own-Topic Project (K--5)\\\large Three-Session Lesson Packet}
\author{(Adapted from teaching\_placement lesson packet conventions)}
\date{March 26, 2026}

% --- Small helpers
\newcommand{\Session}[1]{\section*{#1}}
\newcommand{\Sub}[1]{\subsection*{#1}}
\newcolumntype{Y}{>{\raggedright\arraybackslash}X}

\begin{document}
\maketitle

\Session{Quick Scan}
\begin{itemize}[leftmargin=*]
  \item Grades: K--5 (adaptations included for K--2 and 3--5)
  \item Sessions: 3 (45--60 minutes each; may be shortened to 30--40 for K--2)
  \item Core project: student chooses a topic and creates a small ``Learning Showcase'' product (poster / mini-book / short slide deck)
  \item AI platform: MagicSchool / MagicStudent (exact tool availability depends on school configuration)
  \item Key routine: Plan $\rightarrow$ Build in chunks $\rightarrow$ Validate and revise
\end{itemize}

\Session{School Planning Goals and Inclusion}
\begin{itemize}[leftmargin=*]
  \item \textbf{What is expected:}
  Students use AI as a learning helper (not a replacement), follow privacy rules, and produce an observable artifact that matches a plan.
  \item \textbf{Inclusion focus:}
  visual modeling, chunked steps, predictable routines, partner roles, flexible ways to show learning (oral, visual, typed, dictated).
  \item \textbf{Known friction points to guard against:}
  tool navigation overload, typing/spelling load, and ``decorating'' before meeting success criteria.
  \item \textbf{Teacher success criteria (visible in the room):}
  (1) Everyone has a topic + plan, (2) Everyone builds in chunks using the plan, (3) Everyone runs a final check and can explain what they did.
\end{itemize}

\Session{Digital Citizenship and Safe AI Rules (Post as an Anchor Chart)}
\begin{itemize}[leftmargin=*]
  \item \textbf{No personal information:} do not type full names, addresses, phone numbers, passwords, or private stories.
  \item \textbf{AI can be wrong:} check important facts using class sources or teacher help.
  \item \textbf{AI is a helper:} you still do the thinking and decide what to use.
  \item \textbf{Be kind and school-appropriate:} prompts must match classroom expectations.
  \item \textbf{Teacher visibility:} remind students that classroom AI work is monitored for safety and learning.
\end{itemize}

\Session{Session Overview Table}
\begin{center}
\begin{tabularx}{\textwidth}{@{}p{2.8cm}YYY@{}}
\toprule
\textbf{Session} & \textbf{Objectives} & \textbf{Activities} & \textbf{Assessment / Evidence} \\
\midrule
Session 1: Brainstorm + Plan &
Choose topic; generate 3--5 questions/subtopics; create a plan and success checklist. &
Mini-lesson + modeling; topic narrowing; complete plan template; teacher approval. &
Plan check (topic + questions + steps); exit ticket: ``My topic / my 3 questions / next step''. \\
\addlinespace
Session 2: Build in Chunks &
Draft section-by-section using plan; verify key facts; begin final product. &
Chunked work blocks (Section A/B/C); mid-lesson stop-and-check; partner read-back. &
Teacher observation checklist + 1 screenshot/photo of work-in-progress + oral explanation. \\
\addlinespace
Session 3: Revise + Validate &
Finalize product; use AI to check structure vs plan; revise and reflect. &
AI validation prompt; peer feedback; revise; share; reflect. &
Final product + reflection (how AI helped, what I fixed, what I learned). \\
\bottomrule
\end{tabularx}
\end{center}

% -------------------------
\Session{Session 1: Brainstorm and Lock Topic + High-Level Plan}
\Sub{Teacher materials}
Devices; projector; planning template (below); optional topic bins (books, pictures, short texts); sticky notes; timer.

\Sub{Timing (suggested 50 minutes)}
\begin{itemize}[leftmargin=*]
  \item 0--8 min: Hook + norms (Safe AI Rules)
  \item 8--18 min: Teacher model: ``Bad topic vs good topic'' + how to ask AI for ideas safely
  \item 18--38 min: Student work time: topic choice + plan template
  \item 38--46 min: Quick conferences: approve topic + plan OR revise
  \item 46--50 min: Exit ticket + share 1 idea
\end{itemize}

\Sub{Teacher notes (step-by-step)}
\begin{enumerate}[leftmargin=*]
  \item \textbf{Explain the goal in kid language:}
  ``You will pick something you care about, learn about it, and teach us using a small project. AI can help us think, but we stay in charge.''
  \item \textbf{Model the first screen/tool move (before release):}
  Show exactly where students click to enter the student tool you want them to use (e.g., Idea Generator).
  \item \textbf{Model a safe, age-appropriate prompt:}
  Use a class example topic. Demonstrate asking for topic ideas plus questions.
  \item \textbf{Name the success criteria:}
  Topic + 3--5 questions + product choice + checklist of steps.
  \item \textbf{Conference for approval:}
  Students must show the plan before using AI for drafting next session.
\end{enumerate}

\Sub{Student Plan Template (print or copy into a doc)}
\begin{itemize}[leftmargin=*]
  \item My topic is: \underline{\hspace{10cm}}
  \item Why I chose it (1 sentence or drawing): \underline{\hspace{10cm}}
  \item My 3--5 questions / parts:
    \begin{enumerate}
      \item \underline{\hspace{9cm}}
      \item \underline{\hspace{9cm}}
      \item \underline{\hspace{9cm}}
      \item \underline{\hspace{9cm}}
      \item \underline{\hspace{9cm}}
    \end{enumerate}
  \item My product choice (circle one): Poster \quad Mini-book \quad Slides
  \item My checklist (what I will finish):
    \begin{enumerate}
      \item Plan approved
      \item Section 1 drafted
      \item Section 2 drafted
      \item Section 3 drafted
      \item Visuals added
      \item Final check + revise
      \item Share + reflect
    \end{enumerate}
\end{itemize}

\Sub{Sample student prompts (MagicSchool / AI)}
\begin{itemize}[leftmargin=*]
  \item \textbf{K--2 prompt (Idea Generator):}
  ``I am a student. Give me 5 topic ideas about \underline{\hspace{2cm}} that are easy for kids. For each topic, give 2 questions I could learn about.''
  \item \textbf{3--5 prompt (Idea Generator):}
  ``My interests are \underline{\hspace{2cm}}. Help me choose a focused topic I can teach in 3 sections. Give me 5 options and 4 research questions for each.''
  \item \textbf{Topic narrowing prompt:}
  ``My topic is too big: \underline{\hspace{2cm}}. Help me make it smaller so I can teach it in 3 sections (A, B, C). Give 3 smaller choices.''
\end{itemize}

% -------------------------
\newpage
\Session{Session 2: Learn and Execute Step-by-Step Using the Plan}
\Sub{Teacher materials}
Student plans; note-catcher; teacher-provided sources (books/articles/videos pre-approved); creation tool (poster paper, slides, book tool); timer.

\Sub{Timing (suggested 55 minutes)}
\begin{itemize}[leftmargin=*]
  \item 0--7 min: Re-launch norms + show ``Plan $\rightarrow$ Build in chunks''
  \item 7--12 min: Model: how to draft \emph{one} section from a plan question
  \item 12--42 min: Chunked work time (3 chunks of ~10 min + quick resets)
  \item 42--48 min: Midpoint stop-and-check (compare to plan)
  \item 48--55 min: Save work + quick share (1 fact + how you checked it)
\end{itemize}

\Sub{Teacher notes (step-by-step)}
\begin{enumerate}[leftmargin=*]
  \item \textbf{Start with the plan:}
  Students point to their first question/section before opening tools.
  \item \textbf{Limit tool choice:}
  Preselect 1--2 student tools for drafting (e.g., Informational Texts OR Research Assistant).
  \item \textbf{Teach verification simply:}
  ``If AI says a fact, find it in a class book/text or ask the teacher.''
  \item \textbf{Use a help ladder:}
  Look at directions $\rightarrow$ try 1 small step $\rightarrow$ ask partner $\rightarrow$ ask teacher 1 clear question.
  \item \textbf{Midpoint stop-and-check:}
  Students use a checklist: ``I finished section 1/2/3'' and name what is missing.
\end{enumerate}

\Sub{Sample student prompts (MagicSchool / AI)}
\begin{itemize}[leftmargin=*]
  \item \textbf{K--2 drafting prompt (Informational Texts / Content Creator):}
  ``Write 5 simple sentences about \underline{\hspace{2cm}} for kids. Add 3 vocabulary words and explain them in kid words.''
  \item \textbf{3--5 drafting prompt (Informational Texts):}
  ``Write a short informational passage (3 paragraphs) about \underline{\hspace{2cm}}. Organize it into: Section A \underline{\hspace{2cm}}, Section B \underline{\hspace{2cm}}, Section C \underline{\hspace{2cm}}. Use grade \underline{\hspace{1cm}} words.''
  \item \textbf{Question-by-question prompt (Research Assistant):}
  ``My question is: \underline{\hspace{6cm}}. Give a short answer and 3 key words I can use to check this in a book or class source.''
  \item \textbf{Clarity prompt (Writing Feedback):}
  ``Here is my section. Tell me one thing that is clear and one thing I should fix to make it easier to understand for kids.'' 
\end{itemize}

% -------------------------
\newpage
\Session{Session 3: Finish Project + AI Validation Against Plan \& Lesson Structure}
\Sub{Teacher materials}
Rubric; student plans; peer feedback protocol (two stars + one wish); reflection sheet; sharing setup (gallery walk / small presentations).

\Sub{Timing (suggested 50 minutes)}
\begin{itemize}[leftmargin=*]
  \item 0--6 min: Re-launch: ``Today we polish and check our work.''
  \item 6--18 min: AI validation + structure check vs plan
  \item 18--35 min: Revision time (fix missing parts, improve clarity, correct errors)
  \item 35--45 min: Share (gallery walk or mini-presentations)
  \item 45--50 min: Reflection + submit
\end{itemize}

\Sub{Teacher notes (step-by-step)}
\begin{enumerate}[leftmargin=*]
  \item \textbf{Teach ``AI as checker'':}
  Explain that AI can help check structure/clarity but cannot be trusted automatically for truth.
  \item \textbf{Run the validation prompt:}
  Students paste their plan + their draft into the AI tool (if allowed) and ask for a checklist of missing pieces.
  \item \textbf{Human verification:}
  Require students to confirm at least 2 key facts with a classroom source or teacher check.
  \item \textbf{Celebrate learning:}
  Students share 1 thing they learned and 1 way they used AI responsibly.
\end{enumerate}

\Sub{Sample AI validation prompts}
\begin{itemize}[leftmargin=*]
  \item ``Here is my plan: (paste plan). Here is my project: (paste text). Make a checklist: What matches? What is missing? What should I fix first?''
  \item ``Does my project have these parts: Title, Section A, Section B, Section C, a picture/visual, and a final reflection sentence? If not, tell me what to add.''
  \item ``Find any sentence that sounds too hard for a \underline{\hspace{1cm}} grader. Suggest simpler words.''
\end{itemize}

% -------------------------
\Session{Assessment}
\Sub{Teacher assessment methods}
\begin{itemize}[leftmargin=*]
  \item \textbf{Formative:} plan approval, observation checklist during chunked drafting, midpoint stop-and-check.
  \item \textbf{Summative:} final product + short reflection + responsible AI use.
  \item \textbf{Evidence artifacts:} 1 saved plan, 1 work-in-progress screenshot/photo, final product.
\end{itemize}

\Sub{Optional student reflection (collect at end)}
\begin{itemize}[leftmargin=*]
  \item My topic was: \underline{\hspace{10cm}}
  \item AI helped me by: \underline{\hspace{10cm}}
  \item I checked my work by: \underline{\hspace{10cm}}
  \item One thing I fixed after checking: \underline{\hspace{10cm}}
  \item I am proud of: \underline{\hspace{10cm}}
\end{itemize}

\newpage
\Session{Official References (Teacher)}
\begin{itemize}[leftmargin=*]
  \item CSTA K--12 CS Standards (Revised 2017): \url{https://csteachers.org/k12standards/}
  \item MagicSchool Student Data Policy: \url{https://www.magicschool.ai/privacy-security/student-data-policy}
  \item U.S. Department of Education (OET), AI and the Future of Teaching and Learning (2023): \url{https://www2.ed.gov/documents/ai-report/ai-report.pdf}
  \item UNESCO Guidance for Generative AI in Education and Research (2023): \url{https://www.unesco.org/en/articles/guidance-generative-ai-education-and-research}
  \item CAST UDL Guidelines: \url{https://udlguidelines.cast.org/}
\end{itemize}

\end{document}
